核心公式汇总见表~\ref{tab:formulas}。

\begin{table}
\caption{C--V法测量半导体杂质分布——核心公式汇总}
\label{tab:formulas}
\begin{ruledtabular}
\begin{tabular}{cccc}
\hline
\multicolumn{1}{c}{\textbf{公式}} & \textbf{物理意义} & \textbf{适用条件} & \textbf{推导来源} \\
\hline

$w = \left[\dfrac{2\epsilon\epsilon_0}{qN_{\mathrm{D}}}(V_{\mathrm{D}} + V_{\mathrm{R}})\right]^{1/2}$
& 空间电荷区宽度与偏压关系
& 单边突变结，均匀掺杂，耗尽近似
& \S1.1 \\
\hline

$C = \dfrac{\epsilon\epsilon_0 A}{w}$
& pn结势垒电容（微分电容）
& 单边突变结，均匀掺杂
& \S1.1 \\
\hline

$\dfrac{1}{C^{2}} = \dfrac{2}{A^{2}q\epsilon\epsilon_0 N_{\mathrm{D}}}(V_{\mathrm{R}} + V_{\mathrm{D}})$
& $1/C^{2}$--$V_{\mathrm{R}}$ 线性关系
& 单边突变结，均匀掺杂
& \S1.1 \\
\hline

$N(w) = \dfrac{2}{q\epsilon\epsilon_0 A^{2}\,\mathrm{d}(1/C^{2})/\mathrm{d}V_{\mathrm{R}}}$
& 任意掺杂分布的杂质浓度
& 单边突变结，耗尽近似
& \S1.1 \\
\hline

$N(w) = \dfrac{C^{3}}{q\epsilon\epsilon_0 A^{2}\,|\mathrm{d}C/\mathrm{d}V_{\mathrm{R}}|}$
& 杂质浓度(利用$C$--$V$斜率)
& 单边突变结，耗尽近似
& \S1.1 \\
\hline

$w = \dfrac{\epsilon\epsilon_0 A}{C}$
& 空间电荷区宽度与电容关系
& 普遍成立（耗尽近似）
& \S1.1 \\
\hline

$v_{\mathrm{i}} \approx \dfrac{C_{\mathrm{x}}}{C_{0}}\,v(t)$
& 测量线路输出与待测电容关系
& $C_{0} \gg C_{\mathrm{x}}$
& \S1.2 \\
\hline

$\Delta f_{\mathrm{N}} = \dfrac{1}{4RC}$
& LPF等效噪声带宽
& 一阶RC低通滤波器
& \S1.3 \\
\hline

$\dfrac{\text{SNR}_{\text{out}}}{\text{SNR}_{\text{in}}} = \sqrt{\dfrac{\Delta f_{\text{in}}}{\Delta f_{\text{out}}}}$
& 锁定放大器信噪比改善
& 白噪声假设
& \S1.3 \\
\hline
\end{tabular}
\end{ruledtabular}
\end{table}

\subsection*{数据处理用公式}

\paragraph{杂质浓度 $N(w)$ 的数值计算}

由离散 $C$--$V_{\mathrm{R}}$ 数据，采用中心差分求导：
\begin{equation}
\left.\frac{\mathrm{d}(1/C^{2})}{\mathrm{d}V_{\mathrm{R}}}\right|_{i}
\approx \frac{(1/C^{2})_{i+1} - (1/C^{2})_{i-1}}{V_{\mathrm{R},i+1} - V_{\mathrm{R},i-1}}
\end{equation}

则：
\begin{equation}
N(w_{i}) = \frac{2}{q\epsilon\epsilon_0 A^{2}}
\cdot \frac{V_{\mathrm{R},i+1} - V_{\mathrm{R},i-1}}{(1/C^{2})_{i+1} - (1/C^{2})_{i-1}}
\end{equation}

其中 $w_{i} = \epsilon\epsilon_0 A / C_{i}$。

\paragraph{均匀掺杂段 $N_{\mathrm{D}}$ 的线性拟合}

对 $1/C^{2}$--$V_{\mathrm{R}}$ 曲线的直线段做线性回归 $y = kV_{\mathrm{R}} + b$：
\begin{equation}
N_{\mathrm{D}} = \frac{2}{A^{2}q\epsilon\epsilon_0 \cdot k},\qquad
V_{\mathrm{D}} = -\frac{b}{k}
\end{equation}

\paragraph{已知常数}
\begin{align}
q &= 1.602 \times 10^{-19}\,\mathrm{C} \\
\epsilon_0 &= 8.854 \times 10^{-14}\,\mathrm{F/cm} \\
\epsilon &= 11.8 \quad (\text{硅}) \\
\epsilon\epsilon_0 &= 1.045 \times 10^{-12}\,\mathrm{F/cm} \\
A &= 5.03 \times 10^{-3}\,\mathrm{cm}^{2} \quad (\text{本实验})
\end{align}
